\section{Data and parameter estimation}\label{sec:p2-methods}
% P:p2-methods-01
Loading is $\alpha=n_{\COtwo,\mathrm{added}}/n_{\mathrm{amine,total}}$, per mole of amine molecules, including for PZ.
Unloaded mass fraction is $w_A=m_A/(m_A+m_w)$; solute-free stock compositions are retained for pressure calculations.
True liquid mole fractions use the total number of molecular and ionic species, whereas PZ nuclear magnetic resonance (NMR) group fractions use total PZ as denominator.
The protonated-carbamate group and dicarbamate group are therefore not directly comparable to DEA true mole fractions.
Total pressure and \COtwo partial pressure are separate observables.

% Generator: docs/scientific/latex/scripts/part2_displays.py; no model evaluation.
% Input: analyses/amp_transfer/inputs/kundu/current-declarations.toml SHA256 3085081adcf6b1d5e797a49ae6b96037cc1b5bc5191f45cff699d8546d774c55
% Input: analyses/amp_transfer/inputs/redesign/refit-306-floor-coordinates.json SHA256 b8671b1abbca5c7a406147dfad2359f1c1666887fab0b4eb67e489fc80abbe99
% Input: analyses/amp_transfer/inputs/redesign/refit-306-floor-endpoint.json SHA256 6d26e1e15c89217c7e29ce1d0772aa44285158071fdf4cc1bd0f6e7fb8fb2107
% Input: analyses/amp_transfer/notebook.qmd SHA256 a46134ecc164b9304778c6d809851a64a128dd91f932e6b898da3ce04bc9fd62
% Input: analyses/amp_transfer/results/kundu/summary.json SHA256 814f2ce5a28047c470bbcab59aadf13984e7ae034b55a4239b8e5b3058d0b55e
% Input: analyses/amp_transfer/results/refit-306/kundu-pressures.csv SHA256 561d08cf41e18d6ac277ee5e2dd207904a8885274d8d45739a0ec3f56d851cdb
% Input: analyses/amp_transfer/results/refit-306/refit-306-floor-fit.json SHA256 d7ee97cb1550d2a6969de406f09048dd4326f025a79523b7e205a5c65439233c
% Input: analyses/amp_transfer/results/refit-306/summary.json SHA256 3912f54eb7821b08e05f902d3001493b269355bea383439cada4e33fc4b4128f
% Input: analyses/amp_transfer/results/refit-306/uncertainty/prediction-bands.csv SHA256 d65f885c32fc9b391305b7e0ac76ec9f0112435455d027391c66dcd1ba7a9ff0
% Input: analyses/amp_transfer/results/refit-306/uncertainty/ranking.csv SHA256 ce80d211f43c53c9f654ec85ebdd5a91ac779a827c5e992415954cacffcab2a6
% Input: analyses/dea_transfer/inputs/runs/B-2-coordinates.json SHA256 7402f93dd3334eae251b38df951d2fdc310a730bf0045871bb83482524e75beb
% Input: analyses/dea_transfer/inputs/runs/B-2-endpoint-declarations.toml SHA256 a0e6cd9a5a7796d892130fc71f756480465898fa28cb0fc8cd371e9f4b730b72
% Input: analyses/dea_transfer/inputs/runs/B-2-endpoint-document.json SHA256 91c2384ef9f7b2a035962aafb72d97435a18a94595c5f16c64c9e2ac18785cdd
% Input: analyses/dea_transfer/notebook.qmd SHA256 c08a11c9e055d810c4b04c5c421a3ac89cf9edea895c7eb2ce3225e22368373d
% Input: analyses/dea_transfer/results/figure-values.csv SHA256 58cbb7b6e6f825b0f4fba1ef19f5393ad469e8aa42a58ece391f3b3d7af4bf2c
% Input: analyses/dea_transfer/results/fits.json SHA256 b936638b0b1f83802809bea15c003c1bbcc229b23824ce61d48ac4123b80af48
% Input: analyses/dea_transfer/results/score.json SHA256 12aa281fa03ed5d34cabbe3a09174b49da105199b15f1b16087ffa33018fd64e
% Input: analyses/dea_transfer/results/uncertainty-307/prediction-bands.csv SHA256 b59d9515e2d86b3c76eea92e97343eb5b2e28512bf4e8dd9f69e359de593f9ed
% Input: analyses/dea_transfer/results/uncertainty-307/ranking.csv SHA256 14ec3184971afe521ae4be0f1d636acd3a47067df1a01feb6068d75cf63cf155
% Input: analyses/mdea_mea_workflow/notebook.qmd SHA256 55831c2a23404765f14fbcbeb2c83eb64b7245a10730ad9ad9329b85aa1c2533
% Input: analyses/mdea_mea_workflow/results/refit-298/B1-endpoint/endpoint.json SHA256 2a5341be176829d087133a1467828f762a655f92a586eb1062a5719a58142d9d
% Input: analyses/mdea_mea_workflow/results/refit-298/B1-endpoint/fitted-replay.json SHA256 6690d0b7baf4bb3f08a9437743857d965d7a4a4fd1529d4e6f8e5d234f4059e2
% Input: analyses/mdea_mea_workflow/results/refit-298/attempts.json SHA256 40e17071831a699fe380d4ac4b1b286d5d283ad434a0e704ac2803d7d22c954f
% Input: analyses/mdea_mea_workflow/results/refit-298/pressure-comparison.csv SHA256 d2cdaf4a3bc1a2b7b9b957585a6dbaec2eff15ca56c009b862dac63a2912c22c
% Input: analyses/mea_parameter_bundle/results/selected-current-best-parameters.json SHA256 f8f64820b161377a0d97d1e4ad4f6cea764cd0558dc63231d58e2096988eef97
% Input: analyses/pz_transfer/notebook.qmd SHA256 a71c49071169b4a2795c538b22d2e59c58870489e0131cecf480d5fc98622858
% Input: analyses/pz_transfer/results/holdout-A12.csv SHA256 89f6aabc166a2663ddf7be4cfdc20821cc3ef9ef46070535644f5cd2abb66b15
% Input: analyses/pz_transfer/results/reactive-A12-record.json SHA256 3eefe76e4599050e524b9a1a8354121caae0ca30de38613359cf7e1f9edc1329
% Input: analyses/pz_transfer/results/reactive-A12-residuals.csv SHA256 3c8855bba6dfaa8ed8c8e64f7c4a160cafb77fb1d85b6858e8d5520e1469e4e4
% Input: analyses/pz_transfer/results/reactive-A12-start2-endpoint.json SHA256 c24ff09badc642d5dd312f5dc03248f2e5fa4c93e3b1d91fa69fe42120297b9e
% Input: analyses/pz_transfer/results/speciation-reactions.json SHA256 64a51633b7d56ce75a1246fe51eb6a39a1f0db1abc7d75b16b2fe13df9c125be
% Input: analyses/pz_transfer/results/uncertainty-A12/ranking.csv SHA256 f266707b75ee505bb8a5226c79e6f30037696720f859ca7d239972e6f58c71fe
% Input: analyses/solvent_specification_mdea/B1-declarations.toml SHA256 7b5fd083d83731c74624707f6e9f92bad05265c9ed19d105b25cf99dd897a042
% Input: analyses/solvent_specification_mdea/inputs/B1-coordinates.json SHA256 0bca9af202313b2bd8325ccaa37a72f32a1b326099c957c5ebf35f630506c8be
% Input: analyses/solvent_specification_mdea/notebook.qmd SHA256 227db268f1475ce3e11e4873c68e3337bf217918fa0ce0bdc11deeb25f57ff88
% Input: analyses/solvent_specification_mdea/results/adoption-302/B1/pressure.json SHA256 b81bfeb21d211d64211a2bda07b5e07af2a897f828a750f01aac6142340a7820
% Input: analyses/solvent_specification_mdea/results/retained-statistics-320.json SHA256 64393229910bf06ba614d1345e1fd2d2f675e2aba1a3a9a7eaedcd33ffd07e81
% Input: analyses/solvent_specification_mdea/results/uncertainty-307/prediction-bands.csv SHA256 96fc9352b6f30dae38805b799344404528d6e49b3b45eaa42002daa8523ac0cd
% Input: analyses/solvent_specification_mdea/results/uncertainty-307/predictions.json SHA256 8fd82d94b500a2b8985ea9979e4ceb653fbfb781dc66a3d047ee68bdec0f46ba
% Input: analyses/solvent_specification_mdea/results/uncertainty-307/ranking.csv SHA256 16b4ed6165c0cc7bad4d938567c8f8e8d88bdddb3de4eb8b53c43b372f7c8e7e
% Input: docs/scientific/latex/references.bib SHA256 9ac4ac1c25d01550aa5e9a9df066b7a49e96212687fb4c6f9c826ee81cea4252
% Input: docs/scientific/latex/supplement.tex SHA256 b90143f6a0711b9a7437b4845c8a5db430f8b301ff21d8f402821b105da3d5f8
\begingroup\small\setlength{\tabcolsep}{3pt}
\begin{longtable}{@{}>{\raggedright\arraybackslash}p{0.2\linewidth}>{\raggedright\arraybackslash}p{0.48\linewidth}>{\raggedright\arraybackslash}p{0.25\linewidth}@{}}
\caption{Observation roles and prior access. Held out means never in any fitting objective. Comparison-only means calculated after parameter selection without entering that selection. Prior scoring is stated even when no held-out value determined the objective.}\label{tab:p2-data-roles}\\
\toprule Amine / role & Observations & Scored before final fit chosen? \\ \midrule\endfirsthead
\toprule Amine / role & Observations & Scored before final fit chosen? \\ \midrule\endhead
\bottomrule\endfoot
MDEA fitted & 278 entries: 67 Kuranov; 16 Jakobsen; 40 Carson; 4 Merkley; 33 heat capacities; 33 loaded and 24 unloaded densities; 61 Kim binary pressures & --- \\
MDEA held-out & 20 Xiao and 88 Ermatchkov partial pressures & Yes; previously accessed in model development \\
MDEA exclusions & Kuranov source rows 28 and 72; seven states above the 393.15 K R2 limit & Not counted as available predictions \\
AMP fitted & 156 entries: 47 Hartono loaded total pressures; 52 Hartono binary entries; 57 Ciftja fractions & --- \\
AMP held-out & 51 Kundu partial pressures & Yes; scored during the sequence preceding final diameter choice \\
AMP held-out comparison & 17 Tong total pressures scored; 23 not calculated & Earlier fits had been compared; no blind test claimed \\
PZ fitted & 294 entries: 62 Hilliard pressures; 232 NMR group observations at 121 states & --- \\
PZ held-out & 52 Ermatchkov partial pressures & Yes; unconstrained refit scored before penalty alternative and final choice \\
DEA fitted & 80 entries: 30 Lawson pressures at 25 wt\%; 50 Carson heats & --- \\
DEA held-out & 16 Seo and Hong; 44 Lal partial pressures & Yes; first selected start, then numerically tied second start \\
DEA comparison-only & Lawson Table VIII; Cai binary; Böttinger species; Horstmann excess heat; Weiland densities & No; added after fit selection \\
\end{longtable}\endgroup


% P:p2-methods-02
MDEA uses 278 objective rows: Kuranov et al. total pressures \cite{kuranovSolubilitySingleGases1996}, Jakobsen et al. protonation fractions \cite{Jakobsen2005}, Carson et al. \cite{carsonEnthalpySolutionCarbon2000} and Merkley et al. \cite{merkleyEnthalpiesAbsorptionCarbon1987} heats, Weiland et al. heat capacities \cite{weilandHeatCapacityAqueous1997} and loaded densities \cite{Weiland1998}, Sobrino et al. unloaded densities \cite{sobrinoViscosityDensityMeasurements2016} and Kim et al. binary pressures \cite{kimEbulliometricDeterminationVaporLiquid2008}.
Its Xiao \cite{xiaoCO2AbsorptionAqueous2018a} and Ermatchkov \cite{ermatchkovSolubilityCarbonDioxide2006} pressure sets had been accessed during earlier development and are held-out assessments rather than new blind tests.
AMP uses 47 loaded total pressures \cite{hartonoNewSolubilityHeat2021}, 52 binary observations \cite{hartonoBinaryTernaryVLE2013} and 57 carbamate fractions \cite{ciftjaExperimentalStudyCarbamate2014a}; its Kundu set \cite{kunduVaporLiquidEquilibriumCO22003} was scored before the final diameter restriction was chosen.
PZ uses 62 Hilliard partial pressures \cite{Hilliard2008} and 232 NMR group observations from 121 source states \cite{ermatchkovChemicalEquilibriumConstants2003}.
Its Ermatchkov pressure set was scored on the unconstrained refit before the penalty alternative was examined and the reported fit was finalized.
DEA uses 30 Lawson partial pressures \cite{lawsonGasSweeteningData1976} at 25 wt\% and 50 Carson integral heats at 298.15 K; its held-out sets were first scored on one start and then on a numerically tied second start.
Thus held-out here has the same objective-exclusion meaning for every amine, with prior access reported explicitly.

% P:p2-methods-03
Carrier gases are omitted from the bubble calculations: the model compares $y_{\COtwo}P$ with measured partial pressure rather than imposing a nitrogen- or methane-containing experimental gas.
This affects the MDEA and PZ nitrogen qualifications and the unreported methane pressure in the Lawson DEA cell.
Lal's DEA stock, 2.0 kmol/m$^3$, is represented by mass fraction 0.2058 using unloaded-solution densities.
B\"ottinger DEA comparisons \cite{Bottinger2008} use 2.5 MPa because source cell pressures span 5--25 bar without individual values.
PZ NMR solvent masses reported for heavy water are treated as water, leaving isotope effects unquantified.
These assumptions are part of the compared states, not corrections inferred from the pressure errors.

\subsection{Caloric and density calculations}
% P:p2-methods-05-caloric
A negative pressure bias means underprediction.
For negative absorption heats, the signed difference $y^c-y^o$ instead states directly whether absorption is more or less exothermic.
Supplementary Section~S2 gives the caloric and density definitions, including the stock and loading bases used in each comparison.

\subsection{Parameters and fitting}
% P:p2-methods-04
Each bounded least-squares fit minimizes its original weighted sum,
\begin{equation}
 C=\frac12\boldsymbol r^{\mathsf T}W\boldsymbol r,
 \qquad r_j=(y_j^c-y_j^o)/s_j\quad\text{or}\quad r_j=\ln(y_j^c/y_j^o)/s_j .
 \label{eq:p2-objective}
\end{equation}
Here $s_j$ denotes the assigned observation scale and $W$ includes group weights or correlations.
The objectives use different observations and scales, so their costs do not measure relative physical accuracy across amines.
For DEA, pressure log ratios use scale 0.09 and heat differences use 0.02 times the magnitude of measured heat, with each group's squared residuals averaged.
For PZ, pressure uses an operational 2\,\% scale and NMR uses correlated area-normalization errors with assumed area scale $q=0.02$.
Supplementary Section~S2 specifies the other residuals and their data definitions.

% P:p2-methods-05
The pressure AARD definition is identical to Part I,
\begin{equation}
 \operatorname{AARD}=\frac{100\,\%}{n}\sum_{j=1}^n\left|y_j^c/y_j^o-1\right|,
 \qquad \operatorname{bias}=\frac{100\,\%}{n}\sum_{j=1}^n(y_j^c/y_j^o-1).
 \label{eq:p2-aard}
\end{equation}
The same reporting formula is applied separately to nonzero heats and other properties; species differences and root-mean-square errors retain their observable units.

% P:p2-methods-06
The AMP results use the fit whose carbamate diameter is bounded below by the MEA carbamate diameter, 5.32~\AA; a fit allowed down to 3.5~\AA\ lowers the cost only slightly, so the data identify the diameter only weakly below about 5.3~\AA, and that fit is a sensitivity case.
The PZ fit stops at its 100-iteration limit with nearly flat cost and is not labeled formally converged.
DEA's two charged-chain starts differ in cost below the optimizer tolerance and represent a numerical tie rather than distinct physical optima.
The selected MDEA fit was chosen partly for its held-out pressures and gives poorer agreement with Carson's heats than the earlier fit (scaled RMS 1.106 against 0.977).
Supplementary Section~S3 reports these selections and their parameter identities.

% P:p2-methods-08
Local covariance holds shared inputs, model choices, residual weights and bound coordinates fixed.
For interior coordinates, the covariance is the local inverse weighted information matrix multiplied by the retained residual-variance factor.
At a prediction with parameter derivative $g$, the conditional variance is $g^{\mathsf T}C_\theta g$; displayed intervals use $\pm1.96$ conditional standard errors (SE), not total prediction intervals.
Separate sensitivities displace one marginal SE at a time and do not include correlation compensation.
Inward profiles fix a bound coordinate 1, 5 or 10\,\% of its original interval inside the bound and refit the remaining coordinates.
Those displacements have different meanings and sizes; none quantifies source systematics or model-form error.

\subsection{Phase-stability search}
% P:p2-methods-07
Reacting charge-neutral tangent-plane searches follow Part I, requiring a negative distance and a material-conserving Gibbs-energy decrease to identify an additional phase \cite[Sect.~3]{PolleyPartI}.
Unlike the penalized MEA fit of Part I \cite[Sect.~4.5]{PolleyPartI}, the reported MDEA, AMP, PZ and DEA fits contain no stability penalty; their phases were searched after fitting.
A separate PZ fit applies that penalty and is reported only to show its remaining false liquid splits and the finite search's missed solutions.
All split counts are lower bounds; absence of a detected split is not proof of stability.
The sampled states and unavailable states remain distinct from a continuous phase boundary.
